\documentclass{article}
\usepackage[utf8]{inputenc}
\usepackage{geometry}
\usepackage{graphicx}
\usepackage{pgfplots}
\usepackage{amsmath}
\usepackage{amssymb}
\usepackage{amsfonts}
\usepackage{mdframed}
\usepackage{amsthm}
\usepackage{physics}
\usepackage{derivative}
\usepackage{hyperref}
\usepackage{wrapfig}
\renewcommand{\theenumi}{\roman{enumi}}

% \geometry{a4paper, left=2cm, right=2cm, top=1cm, bottom=2cm}

\pgfplotsset{compat=1.18}

\date{10 de diciembre del 2024}
\begin{document}
\input{frontpage}



\title{Fenomenología de Colisiones Galácticas: Efectos Extragalácticos y el Rol de la Materia Oscura en la Dinámica Universal}
\maketitle

\section*{Introducción}
Las colisiones galácticas representan uno de los procesos más espectaculares y significativos en la evolución del universo. A lo largo del tiempo cósmico, las galaxias, al estar inmersas en un entorno gravitacional dinámico, interactúan y colisionan en encuentros que pueden abarcar desde perturbaciones menores hasta fusiones completas. Estos eventos, regidos principalmente por las fuerzas gravitacionales, son facilitados por la estructura jerárquica del universo, en la cual las galaxias y los cúmulos de galaxias tienden a agruparse bajo la influencia de halos masivos de materia oscura. \\

El concepto de colisión galáctica no implica necesariamente choques directos entre estrellas debido a las enormes distancias que las separan, sino más bien interacciones gravitacionales que redistribuyen la materia estelar, el gas y el polvo. Sin embargo, estas interacciones generan dinámicas complejas, como ondas de choque en el gas interestelar, fenómenos de marea gravitacional y formación de nuevas estrellas debido a la compresión del medio interestelar. Ejemplos icónicos, como la interacción entre las galaxias de Las Antenas, ilustran cómo estas colisiones pueden reconfigurar de manera dramática la morfología galáctica, creando puentes de materia, colas de marea y regiones intensas de formación estelar. \\

En términos cosmológicos, las colisiones galácticas son eventos inevitables en el universo en expansión. Se estima que la mayoría de las galaxias grandes, incluida nuestra Vía Láctea, han experimentado múltiples colisiones a lo largo de su historia. La próxima fusión entre la Vía Láctea y la galaxia de Andrómeda, predicha para ocurrir dentro de unos 4 mil millones de años, representa un caso emblemático que subraya la relevancia de estos eventos en la evolución futura del universo local. \\

El estudio de las colisiones galácticas es un componente esencial de la astrofísica extragaláctica, ya que estos eventos ofrecen una ventana única para investigar los procesos que moldean la evolución de las galaxias y sus entornos. Las interacciones galácticas permiten a los astrofísicos explorar fenómenos físicos fundamentales, como la transferencia de momento angular, la redistribución de masa, y la formación de estructuras a gran escala, como cúmulos y filamentos galácticos. Asimismo, proporcionan información crítica sobre los ciclos de vida del gas interestelar y la formación de estrellas, que son procesos clave para entender cómo las galaxias transforman su contenido químico y energético a lo largo del tiempo. \\

En términos observacionales, las colisiones galácticas ofrecen pistas cruciales sobre la distribución y propiedades de la materia oscura, un componente invisible pero dominante en la dinámica galáctica. Las simulaciones numéricas han mostrado que los efectos gravitacionales observados durante las colisiones no pueden explicarse únicamente por la materia visible, lo que refuerza la hipótesis de que halos masivos de materia oscura envuelven a las galaxias. Este vínculo directo convierte a las colisiones galácticas en un laboratorio natural para probar teorías cosmológicas y modelos de materia oscura. \\

Además, las interacciones galácticas tienen implicaciones significativas para la evolución del universo a gran escala. A medida que las galaxias colisionan y se fusionan, contribuyen a la formación jerárquica de estructuras mayores, desde galaxias elípticas masivas hasta cúmulos y supercúmulos. Este marco jerárquico no solo respalda el modelo cosmológico estándar basado en la materia oscura fría y la energía oscura, sino que también plantea preguntas abiertas sobre la dinámica interna de los sistemas galácticos y su evolución futura. \\

Por estas razones, el estudio de las colisiones galácticas no solo enriquece nuestra comprensión de las galaxias como sistemas individuales, sino que también conecta estos procesos con el marco cosmológico más amplio, proporcionando una visión integral del universo y sus componentes fundamentales. \\

La materia oscura desempeña un papel central en la dinámica de las colisiones galácticas y en la evolución del universo en general. Aunque su naturaleza exacta permanece desconocida, se estima que constituye aproximadamente el 27\% del contenido total del universo, superando ampliamente la materia ordinaria. Este componente invisible, detectable únicamente a través de sus efectos gravitacionales, forma halos masivos alrededor de las galaxias, los cuales actúan como una "estructura de soporte" que guía y condiciona sus interacciones y movimientos. \\

Durante una colisión galáctica, la materia oscura influye directamente en la trayectoria, duración y resultados de la interacción. A diferencia de las estrellas y el gas, que pueden experimentar fricción y redistribuirse mediante interacciones físicas, la materia oscura se comporta como un fluido colisionalmente débil, interactuando principalmente a través de la gravedad. Este comportamiento único afecta no solo la dinámica de las galaxias en colisión, sino también la formación de estructuras extragalácticas, como puentes y colas de marea, que a menudo se extienden mucho más allá de las regiones visibles. \\

Además, los efectos de la materia oscura pueden observarse en la redistribución del momento angular y en la estabilización de sistemas galácticos fusionados. Por ejemplo, los halos masivos de materia oscura contribuyen a la preservación de ciertas estructuras después de una colisión, lo que permite la formación de galaxias elípticas gigantes o sistemas estelares relajados. Simultáneamente, las interacciones gravitacionales generadas por la materia oscura pueden desencadenar ondas de choque y procesos de formación estelar intensiva, especialmente en regiones de alta densidad de gas. \\

Desde una perspectiva cosmológica, las colisiones galácticas ofrecen una oportunidad única para estudiar la distribución y propiedades de la materia oscura. Observaciones de sistemas como los cúmulos galácticos de la "Bala" han revelado discrepancias entre la distribución de materia luminosa y los efectos gravitacionales detectados, proporcionando evidencia robusta sobre la presencia de materia oscura. Por lo tanto, estudiar estas interacciones no solo ilumina los procesos internos de las galaxias, sino que también contribuye a validar y refinar modelos de la estructura del universo. \\
%objetivos
Este trabajo tiene como propósito central explorar en profundidad la fenomenología de las colisiones galácticas, con un enfoque particular en los efectos extragalácticos que generan y en el rol que desempeña la materia oscura en dichos eventos. Los objetivos específicos se muestran a continuación: \\

\begin{enumerate}
    \item Analizar los mecanismos físicos subyacentes en las colisiones galácticas: Examinar los procesos dinámicos y gravitacionales que afectan la redistribución de masa y energía durante estas interacciones.
    \item Investigar los efectos extragalácticos derivados de las colisiones: Estudiar cómo estos eventos influyen en la formación de estructuras a gran escala, como puentes de marea, cúmulos galácticos y sistemas estelares secundarios.
    \item Explorar el papel de la materia oscura en las colisiones galácticas: Evaluar su impacto en las trayectorias de las galaxias, en los resultados observables de las interacciones y en la dinámica gravitacional global.
    \item Revisar casos de estudio observacionales y simulados: Analizar ejemplos emblemáticos de colisiones galácticas, comparando observaciones astronómicas con simulaciones numéricas avanzadas para validar modelos teóricos.
    \item Discutir las implicaciones cosmológicas de las colisiones galácticas: Relacionar estos eventos con la formación jerárquica de estructuras y con las teorías actuales sobre la evolución del universo y la naturaleza de la materia oscura.
\end{enumerate}
Al abordar estos objetivos, el presente estudio pretende no solo contribuir al entendimiento de las colisiones galácticas en sí, sino también arrojar luz sobre los procesos fundamentales que gobiernan la dinámica del cosmos, con especial atención al enigma persistente de la materia oscura. \\

\section{Capítulo 1: Fundamentos Teóricos}
\subsection{Breve repaso de la dinámica galáctica}
La dinámica galáctica estudia los principios físicos que gobiernan el movimiento de las estrellas, el gas y otras componentes en las galaxias bajo la influencia de la gravedad. Este campo integra herramientas de mecánica celeste, hidrodinámica y teoría gravitacional para describir cómo las galaxias evolucionan, interactúan y mantienen su cohesión estructural en el tiempo cósmico. \\

En el caso de las galaxias espirales, como la Vía Láctea, la distribución de masa visible se organiza en un disco plano donde predominan las órbitas casi circulares de las estrellas. Sin embargo, las observaciones de las curvas de rotación galácticas han demostrado que la velocidad orbital de las estrellas y el gas en las regiones exteriores no disminuye como lo predice la ley de Kepler, sino que permanece constante o incluso aumenta. Este fenómeno, inconsistente con la distribución de masa visible, se atribuye a la presencia de halos masivos de materia oscura que extienden su influencia gravitacional más allá de los límites ópticamente visibles de las galaxias. \\

Las galaxias elípticas, por otro lado, exhiben una estructura más tridimensional, con trayectorias estelares que siguen patrones aleatorios y anisotrópicos. En este caso, la dinámica está dominada por la presión estelar y los campos gravitacionales colectivos, que dan lugar a distribuciones de masa esféricas o ligeramente elipsoidales. La falta de rotación global significativa en estas galaxias contrasta con la naturaleza ordenada de las espirales y refleja diferencias en sus historias de formación, muchas veces marcadas por colisiones y fusiones con otras galaxias. \\

En sistemas galácticos en interacción, las fuerzas gravitacionales mutuas generan fenómenos de marea que redistribuyen la masa y el momento angular. Estas fuerzas de marea pueden arrancar estrellas y gas de las galaxias, formando estructuras transitorias como puentes y colas de marea. Por ejemplo, en sistemas interactuantes como Las Antenas, estas dinámicas conducen a una intensa formación estelar en regiones comprimidas por las interacciones gravitacionales. \\

A escala cosmológica, las galaxias no son sistemas aislados, sino que forman parte de estructuras más grandes, como cúmulos y filamentos, vinculados gravitacionalmente dentro de la red cósmica. La interacción de las galaxias con estas estructuras influye en su evolución, afectando la captura de gas, la formación estelar y las colisiones con otras galaxias. La presencia ubicua de materia oscura en estas escalas amplifica los efectos gravitacionales, conectando la dinámica galáctica local con la evolución del universo en su conjunto. \\

Desde el punto de vista teórico, la dinámica galáctica se modela a través de simulaciones numéricas que resuelven las ecuaciones gravitacionales para un gran número de partículas. Estas simulaciones, como las del Proyecto Illustris o Millennium, han sido fundamentales para entender la formación y evolución galáctica, así como para predecir fenómenos observables en colisiones y fusiones. Además, los avances en los modelos computacionales han permitido incorporar no solo la interacción gravitacional, sino también procesos hidrodinámicos, de formación estelar y de retroalimentación, proporcionando una visión más completa de la evolución galáctica.

\subsection{Naturaleza de las colisiones galácticas (fusión, interacción, etc.)}
Las colisiones galácticas son encuentros gravitacionales entre galaxias que, dependiendo de sus condiciones iniciales y de las propiedades físicas de los sistemas involucrados, pueden manifestarse como interacciones menores o eventos catastróficos que culminan en fusiones completas. Estos procesos no solo son comunes en la historia cósmica, sino que también juegan un papel esencial en la evolución de las galaxias, ya que transforman su morfología, dinámica y composición. \\

\subsubsection{Interacciones menores y perturbaciones gravitacionales}
Las interacciones menores ocurren cuando una galaxia más pequeña, a menudo denominada galaxia satélite, interactúa gravitacionalmente con una galaxia más grande sin que necesariamente se produzca una fusión. Estas interacciones suelen inducir perturbaciones en los discos de gas y estrellas de la galaxia principal, generando fenómenos como barras, espirales o deformaciones asimétricas. Un ejemplo prominente es la interacción de las Nubes de Magallanes con la Vía Láctea, que ha resultado en la formación de una corriente de gas conocida como la Corriente Magallánica. \\

Aunque la galaxia satélite puede perder parte de su masa estelar o de gas debido a fuerzas de marea, estas interacciones no suelen alterar drásticamente la estructura global de la galaxia principal. Sin embargo, en muchos casos, el satélite experimenta una disolución progresiva, siendo finalmente asimilado por la galaxia más grande. Este tipo de eventos contribuye a la acreción jerárquica, un proceso clave en el crecimiento galáctico. \\

\subsubsection{Colisiones mayores y fusiones galácticas}
En los encuentros más intensos, las galaxias participantes tienen masas comparables y colisionan directamente. Estas colisiones son complejas, ya que involucran interacciones gravitacionales entre los componentes visibles e invisibles de ambas galaxias, así como procesos hidrodinámicos en el gas interestelar. Dependiendo de los parámetros orbitales y de la cantidad de momento angular involucrado, estos encuentros pueden dar lugar a fusiones completas o a sistemas galácticos profundamente alterados. \\

Las fusiones galácticas transforman de manera radical las propiedades físicas de las galaxias participantes. Por ejemplo, dos galaxias espirales pueden fusionarse y formar una galaxia elíptica masiva, como se predice en la futura colisión entre la Vía Láctea y Andrómeda. Este tipo de transformación está impulsado por la redistribución del momento angular y la disipación de energía orbital en forma de calor interno, lo que resulta en la mezcla de los discos estelares y en la creación de una estructura más esférica. \\

Durante las fusiones, las fuerzas de marea y las ondas de choque comprimen el gas interestelar, desencadenando episodios intensos de formación estelar. Estos brotes de formación estelar, conocidos como "starbursts", son responsables de producir grandes cantidades de nuevas estrellas en un periodo de tiempo relativamente corto. En casos extremos, la energía liberada por la formación estelar y la actividad de núcleos galácticos activos (AGN) puede expulsar parte del gas galáctico al medio intergaláctico, deteniendo posteriormente la formación de nuevas estrellas. \\

\subsubsection{Duración y escala temporal de las colisiones}
Las colisiones galácticas no son eventos instantáneos, sino procesos que se desarrollan en escalas de tiempo de cientos de millones a miles de millones de años. Las fases iniciales incluyen la aproximación y la generación de fuerzas de marea, seguidas por una etapa de contacto cercano, donde el gas y las estrellas comienzan a redistribuirse. En casos de fusión, las galaxias pueden experimentar múltiples pasos orbitales antes de integrarse completamente en un único sistema galáctico. \\

\subsubsection{Impacto de la materia oscura}
La naturaleza de las colisiones galácticas está profundamente influenciada por la materia oscura. Los halos masivos de materia oscura no solo definen las trayectorias iniciales de las galaxias, sino que también estabilizan y amplifican las interacciones gravitacionales. En simulaciones numéricas, se observa que la fusión de los halos de materia oscura ocurre antes que la mezcla completa de las componentes visibles, lo que subraya su papel dominante en la dinámica general. \\

\subsubsection{Implicaciones cosmológicas}
A nivel cosmológico, las colisiones y fusiones galácticas son fundamentales en el modelo jerárquico de formación de estructuras, donde las galaxias más pequeñas se combinan progresivamente para formar sistemas mayores y más complejos. Este marco teórico, respaldado por simulaciones como Millennium y Illustris, conecta los procesos de colisión con la evolución de la red cósmica, contribuyendo al crecimiento de cúmulos y supercúmulos de galaxias.

\subsection{Introducción a la materia oscura: hipótesis actuales y observaciones clave}

La materia oscura es uno de los componentes más misteriosos y fundamentales del universo. Representa aproximadamente el 27\% de su densidad total, superando ampliamente a la materia ordinaria (aproximadamente 5\%) y desempeñando un papel crucial en la formación y evolución de las estructuras cósmicas. Aunque no interactúa con la radiación electromagnética y, por lo tanto, es invisible para los telescopios tradicionales, su presencia se infiere a través de sus efectos gravitacionales en la materia visible, el gas y la luz. \\

\subsubsection{Hipótesis actuales sobre la naturaleza de la materia oscura}
Partículas masivas de interacción débil (WIMPs): \\
Las WIMPs son una de las hipótesis más estudiadas y plantean que la materia oscura está compuesta por partículas masivas que interactúan débilmente con la materia ordinaria. Estas partículas serían relictos del universo temprano, formadas durante el enfriamiento del Big Bang. Las WIMPs son consistentes con modelos teóricos como la supersimetría, pero hasta ahora no se han detectado experimentalmente. \\

\textit{Axiones} \\
Los axiones son partículas ligeras propuestas como solución a problemas en la cromodinámica cuántica. Podrían explicar la materia oscura como un campo cuántico extendido a gran escala en el universo. Experimentos actuales, como ADMX, buscan detectar señales de axiones a través de interacciones con campos magnéticos. \\

\textit{Materia oscura caliente y tibia} \\
Aunque la hipótesis predominante es la materia oscura fría (CDM), también se han propuesto alternativas como la materia oscura caliente, formada por partículas relativistas como neutrinos, o tibia, compuesta por partículas con propiedades intermedias. Estas variantes afectan la formación de estructuras a pequeña escala y podrían explicar discrepancias entre las simulaciones basadas en CDM y las observaciones. \\

\textit{Teorías modificadas de la gravedad} \\
Algunas hipótesis sugieren que los efectos atribuidos a la materia oscura podrían explicarse mediante modificaciones a la teoría de la gravedad, como la Dinámica Newtoniana Modificada (MOND) o teorías relativistas generalizadas. Sin embargo, estas teorías enfrentan dificultades para explicar fenómenos a gran escala, como las anisotropías del fondo cósmico de microondas. \\

\subsubsection{Observaciones clave que respaldan la existencia de la materia oscura}
Curvas de rotación galácticas: \\
Una de las primeras evidencias directas de la materia oscura provino del estudio de las curvas de rotación de galaxias espirales. Las mediciones mostraron que las velocidades orbitales de las estrellas y el gas en las regiones exteriores permanecen casi constantes, en contraste con lo predicho por la distribución de masa visible. Este comportamiento solo puede explicarse asumiendo la presencia de una masa invisible dominante en forma de un halo esférico de materia oscura. \\

\begin{enumerate}
    \item Lentes gravitacionales: La materia oscura también se detecta a través de efectos de lente gravitacional, donde la luz de objetos distantes se curva debido a la presencia de grandes cantidades de masa. Los mapas de lentes gravitacionales, como los obtenidos en el cúmulo de galaxias Abell 1689, han mostrado que la masa total de estos sistemas excede significativamente la masa visible.
    \item Colisiones de cúmulos de galaxias: Un ejemplo icónico es el cúmulo de la Bala, donde el gas caliente, detectable en rayos X, está desplazado de los centros de masa inferidos mediante lentes gravitacionales. Esto sugiere que la materia oscura no interactúa significativamente con el gas ni consigo misma, reforzando su carácter colisionalmente débil.
    \item Anisotropías del fondo cósmico de microondas (CMB): Las observaciones del CMB, como las realizadas por los telescopios Planck y WMAP, han proporcionado una visión detallada de las fluctuaciones de densidad primordiales en el universo temprano. Los modelos que incluyen materia oscura fría se ajustan con precisión a las propiedades observadas del CMB, lo que refuerza su papel en la formación de estructuras cósmicas.
    \item Formación de estructuras a gran escala: Simulaciones numéricas, como Millennium e Illustris, muestran que la materia oscura es fundamental para el crecimiento jerárquico de estructuras cósmicas. Sin su presencia, las galaxias y cúmulos no podrían formarse en los tiempos observados, ya que la materia ordinaria por sí sola no sería capaz de condensarse con suficiente rapidez debido a la presión térmica.
\end{enumerate}

\subsubsection{Implicaciones de la materia oscura en la astrofísica extragaláctica}
La materia oscura no solo es esencial para comprender la dinámica interna de las galaxias, sino que también actúa como el "pegamento gravitacional" que conecta y estabiliza las estructuras a gran escala del universo. En el contexto de las colisiones galácticas, su influencia es decisiva para modelar trayectorias, redistribuir masa y definir los resultados finales de estas interacciones. Al mismo tiempo, estudiar estas colisiones ofrece una oportunidad única para refinar las hipótesis sobre su naturaleza, aportando datos clave que podrían resolver uno de los mayores misterios de la cosmología moderna.

\section{Capítulo 2: Mecanismos de Colisiones Galácticas}
\subsection{Descripción de las interacciones gravitacionales}
Las interacciones gravitacionales son el motor fundamental detrás de las colisiones galácticas. Estas fuerzas actúan sobre las componentes estelares, gaseosas y de materia oscura de las galaxias en interacción, desencadenando procesos complejos que transforman la estructura, la dinámica y la composición de los sistemas involucrados. Este apartado detalla los principios físicos subyacentes y los fenómenos asociados con estas interacciones.

\subsubsection{Principios básicos de las interacciones gravitacionales}
La gravedad, siendo una fuerza de largo alcance, domina las dinámicas de las galaxias, que pueden separarse por decenas o cientos de miles de años luz. A diferencia de las fuerzas electromagnéticas, la gravedad afecta de manera acumulativa a lo largo del tiempo, permitiendo que incluso pequeñas perturbaciones gravitacionales tengan un impacto significativo en las trayectorias y configuraciones de los sistemas galácticos. \\

Durante una colisión galáctica, la interacción gravitacional afecta a tres componentes principales:

\begin{enumerate}
    \item Estrellas: Aunque las estrellas individuales raramente colisionan directamente debido a las grandes distancias relativas entre ellas, su distribución y movimiento son alterados significativamente por las fuerzas de marea. Esto da lugar a fenómenos como puentes y colas de marea.
    \item Gas interestelar: A diferencia de las estrellas, el gas es altamente colisional y responde tanto a las fuerzas gravitacionales como a las interacciones hidrodinámicas, lo que puede generar ondas de choque y episodios de formación estelar.
    \item Materia oscura: Los halos de materia oscura, aunque invisibles, tienen un efecto dominante en las trayectorias y en la estabilización gravitacional durante y después de la interacción.
\end{enumerate}

\subsubsection{Efectos de marea gravitacional}
Uno de los fenómenos más característicos de las interacciones gravitacionales es la generación de fuerzas de marea. Estas fuerzas surgen debido a diferencias en la intensidad de la atracción gravitacional ejercida sobre distintas partes de una galaxia. En un sistema en colisión, estas diferencias deforman las galaxias, extrayendo material estelar y gaseoso hacia el espacio intergaláctico.

Puentes de marea: Estas estructuras conectan temporalmente las galaxias en interacción y están compuestas por estrellas y gas arrancados de los discos galácticos.
Colas de marea: Material estelar y gaseoso es expulsado a grandes distancias, formando estructuras elongadas que pueden persistir durante miles de millones de años.
Ejemplos emblemáticos de estos fenómenos incluyen las galaxias Antennae y el sistema de las galaxias Mice, donde las fuerzas de marea han creado estructuras extensas y llamativas.

\subsubsection{Redistribución de masa y momento angular}
Durante una interacción, la energía orbital de las galaxias se convierte en energía interna de los sistemas, lo que redistribuye la masa y el momento angular. Este proceso da lugar a:

\begin{enumerate}
    \item Contracción de los núcleos galácticos: El material estelar y gaseoso puede ser arrastrado hacia el centro de las galaxias debido a la pérdida de momento angular, generando cúmulos estelares densos o alimentando agujeros negros supermasivos.
    \item Desorganización de discos estelares: En colisiones mayores, los discos de galaxias espirales pueden perder su estructura, formando galaxias elípticas con una distribución de masa más esférica.
    \item Expulsión de material: Parte del gas y las estrellas puede ser expulsado hacia el medio intergaláctico, contribuyendo al enriquecimiento químico de este medio.
\end{enumerate}

\subsubsection{Escala temporal de las interacciones}
Las colisiones galácticas no ocurren de manera instantánea. La evolución completa de un sistema en interacción puede dividirse en varias etapas:

\begin{enumerate}
    \item Aproximación inicial: Las galaxias se atraen gravitacionalmente y comienzan a experimentar fuerzas de marea. En esta etapa, las trayectorias son dominadas por la energía orbital y los halos de materia oscura comienzan a interactuar.
    \item Primer contacto cercano: Las galaxias alcanzan su distancia mínima, lo que genera las perturbaciones gravitacionales más intensas, como la formación de puentes y colas de marea.
    \item Oscilaciones orbitales: Después del primer encuentro, las galaxias pueden separarse parcialmente antes de volver a interactuar debido a la pérdida de energía orbital.
    \item Fusión final: Si la interacción es suficientemente intensa y no hay suficiente momento angular para mantenerlas separadas, las galaxias se fusionan en un único sistema, a menudo formando una galaxia elíptica o irregular.
\end{enumerate}

\subsubsection{Rol de la materia oscura en las interacciones gravitacionales}
Los halos de materia oscura que rodean las galaxias amplifican las interacciones gravitacionales al extender la influencia gravitacional mucho más allá de la distribución visible de masa. Durante una colisión:

\begin{enumerate}
    \item Los halos de materia oscura interactúan primero, "anclando" las trayectorias iniciales y permitiendo que las galaxias visibles se acerquen lo suficiente para interactuar directamente.
    \item La fusión de los halos ocurre antes que la mezcla de las componentes visibles, estabilizando el sistema resultante después de la colisión.
    \item La inclusión de la materia oscura en simulaciones numéricas es esencial para reproducir las características observadas en colisiones galácticas reales, destacando su papel fundamental en la dinámica gravitacional.
\end{enumerate}

\subsection{Procesos de marea y redistribución de masa estelar}
Los procesos de marea representan uno de los aspectos más característicos e impactantes de las colisiones galácticas. Generados por las diferencias en la fuerza gravitacional ejercida sobre diferentes partes de una galaxia, estos procesos conducen a una significativa redistribución de la masa estelar y gaseosa, alterando profundamente la estructura de las galaxias involucradas. Esta sección explora los mecanismos detrás de estos procesos y sus efectos en la dinámica galáctica.

\subsubsection{Fuerzas de marea: principios básicos}
Las fuerzas de marea surgen de gradientes gravitacionales. En una colisión galáctica, las fuerzas gravitacionales que actúan sobre una galaxia no son uniformes:

Las regiones más cercanas al centro de masa de la galaxia compañera experimentan una atracción gravitacional más fuerte.
Las regiones más alejadas sienten una fuerza gravitacional más débil.
Este diferencial crea tensiones internas que distorsionan la galaxia, arrancando estrellas y gas de su estructura original. Este material puede ser redistribuido en el espacio intergaláctico, formando características prominentes como puentes y colas de marea.

\subsubsection{Estructuras producidas por las fuerzas de marea}
Los procesos de marea generan estructuras distintivas que se observan comúnmente en sistemas galácticos interactuantes:

\subsubsection{Puentes de marea}
Los puentes de marea son estructuras transitorias formadas por material arrancado de los discos galácticos que conecta a las galaxias en interacción. Estas estructuras son ricas en gas y estrellas, y su formación depende de la orientación, velocidad y proximidad de las galaxias. Un ejemplo destacado de puentes de marea se encuentra en el sistema de Las Antenas, donde las galaxias están unidas por un puente de gas y estrellas.

\subsubsection{Colas de marea}
Las colas de marea son flujos alargados de estrellas y gas que son expulsados hacia las afueras de las galaxias debido a las fuerzas de marea. Estas estructuras suelen ser más prominentes en colisiones cercanas, donde la interacción gravitacional es más intensa. Las colas de marea pueden extenderse por cientos de miles de años luz y persisten durante largos períodos de tiempo, proporcionando pistas sobre la historia de las interacciones galácticas. Un ejemplo notable es la cola de marea en el sistema de las Galaxias Ratón (NGC 4676).

\subsubsection{Ondas de densidad y resonancias internas}
Además de los flujos de material hacia el espacio intergaláctico, las fuerzas de marea pueden inducir ondas de densidad en los discos galácticos. Estas ondas de densidad reorganizan la masa estelar dentro de las galaxias, formando barras o intensificando patrones espirales.

\subsubsection{Redistribución de masa estelar y su impacto dinámico}
La redistribución de masa estelar en colisiones galácticas tiene implicaciones dinámicas significativas para las galaxias involucradas:

\subsubsection{Migración de estrellas hacia las regiones externas}
Parte de las estrellas arrancadas por las fuerzas de marea se desplaza hacia el espacio intergaláctico o forma estructuras extendidas alrededor de las galaxias, como envolturas difusas de estrellas. Estas poblaciones estelares pueden permanecer gravitacionalmente ligadas o ser completamente expulsadas.

\subsubsection{Concentración de masa en el núcleo galáctico}
En muchos casos, las fuerzas de marea canalizan masa hacia el centro de las galaxias, aumentando la densidad estelar en sus núcleos. Este fenómeno puede desencadenar la formación de cúmulos estelares compactos o alimentar agujeros negros supermasivos. En sistemas con núcleos activos, este proceso puede intensificar la actividad de núcleos galácticos activos (AGN).

\subsubsection{Destrucción de estructuras galácticas preexistentes}
Las interacciones gravitacionales pueden desorganizar discos estelares, transformando galaxias espirales en sistemas elípticos con distribuciones de masa más homogéneas y esféricas.

\subsubsection{Escala temporal de los procesos de marea}
Los procesos de marea se desarrollan a lo largo de varias etapas durante una colisión galáctica:

\begin{enumerate}
    \item Aproximación inicial:
Las galaxias comienzan a experimentar gradientes gravitacionales a medida que se acercan. Este es el momento en que se generan las primeras deformaciones en los discos galácticos.

    \item Encuentro cercano:
En esta etapa, las fuerzas de marea alcanzan su máximo, arrancando material estelar y formando puentes y colas.

    \item Separación parcial:
Las galaxias pueden separarse temporalmente después de su primer paso cercano, con estructuras de marea persistiendo como evidencia de la interacción.

    \item Fusión final:
Si las galaxias pierden suficiente momento angular, se fusionarán completamente, redistribuyendo sus componentes estelares en una nueva configuración.
\end{enumerate}

\subsubsection{Materia oscura y procesos de marea}
Los halos de materia oscura juegan un papel crucial en los procesos de marea al influir en cómo se distribuye y se redistribuye la masa estelar. Estos amortiguan las fuerzas de marea al proporcionar una masa gravitacional dominante que mantiene unida a la galaxia incluso cuando su estructura visible se desintegra parcialmente. \\
Además la presencia de halos masivos de materia oscura amplía el alcance de las interacciones gravitacionales, permitiendo que las galaxias se influencien mutuamente a distancias mayores y durante períodos de tiempo más largos.

\subsection{Simulaciones numéricas: metodología y resultados generales}
Las simulaciones numéricas son herramientas fundamentales en el estudio de las colisiones galácticas, ya que permiten explorar los complejos procesos dinámicos que no pueden ser observados directamente debido a las enormes escalas de tiempo involucradas. Estas simulaciones, basadas en principios gravitacionales, hidrodinámicos y de dinámica de partículas, han transformado nuestra comprensión de cómo las interacciones galácticas moldean la estructura y evolución del universo.

\subsubsection{Metodología de las simulaciones numéricas}
El diseño de una simulación numérica de colisiones galácticas involucra varios pasos fundamentales:

\subsubsection*{Modelado inicial de las galaxias}

\begin{enumerate}
    \item Componentes visibles:
    Las galaxias se representan mediante distribuciones de partículas que simulan las estrellas, el gas interestelar y otros elementos visibles. Estas partículas obedecen a las leyes de la dinámica gravitacional.
    \item Halo de materia oscura:
Los halos galácticos se modelan con distribuciones de partículas de materia oscura, generalmente usando perfiles de densidad como el perfil de Navarro-Frenk-White (NFW), que describe cómo la densidad disminuye con la distancia al centro galáctico.
\end{enumerate}

\subsubsection*{Configuración de condiciones iniciales}

Las posiciones, velocidades y orientaciones iniciales de las galaxias son parámetros cruciales que determinan el desarrollo de la interacción. Estas condiciones incluyen:
Separación inicial entre las galaxias.
Trayectorias relativas (parámetros orbitales).
Distribución interna de masa y energía en cada galaxia.
\subsubsection*{Resolución y algoritmos computacionales}

\begin{enumerate}
    \item Método N-corpos:
Este método calcula las fuerzas gravitacionales entre todas las partículas de la simulación. Sin embargo, el costo computacional crece rápidamente con el número de partículas, por lo que a menudo se utiliza una variante conocida como árbol jerárquico (tree-code) para reducir el tiempo de cálculo.
    \item Dinámica de fluidos:
El gas interestelar se modela mediante hidrodinámica de malla o métodos de partículas suavizadas (SPH, por sus siglas en inglés) para simular interacciones colisionables, como ondas de choque o compresión.
\end{enumerate}

\subsubsection*{Escalado temporal}
Las simulaciones abarcan miles de millones de años en escalas cósmicas, y los resultados se graban en intervalos discretos para estudiar la evolución de las galaxias paso a paso.

\subsubsection{Resultados generales de las simulaciones numéricas}
Las simulaciones numéricas han revelado patrones universales y fenómenos comunes en las colisiones galácticas, que son consistentes con observaciones astronómicas. Entre los resultados más significativos se encuentran.

\subsubsection*{Formación de estructuras características}

Puentes y colas de marea:
Las simulaciones reproducen la formación de estructuras observadas en sistemas interactuantes, como las galaxias Antennae, validando la interpretación de que estas son causadas por fuerzas de marea.
Remanentes de fusión:
Las fusiones mayores tienden a transformar galaxias espirales en galaxias elípticas. Las simulaciones han sido fundamentales para comprender cómo la redistribución de masa y la pérdida de momento angular conducen a esta transición morfológica.

\subsubsection*{Formación estelar y actividad nuclear}
Las interacciones gravitacionales comprimen el gas interestelar, desencadenando explosiones de formación estelar (starbursts). Las simulaciones también muestran que este gas puede canalizarse hacia el núcleo galáctico, alimentando agujeros negros supermasivos y dando lugar a núcleos galácticos activos (AGN).

\subsubsection*{Papel de la materia oscura}
Los halos de materia oscura estabilizan los sistemas galácticos durante las colisiones y absorben parte del momento angular, permitiendo que las galaxias visibles interactúen más intensamente. Las simulaciones sugieren que las distribuciones de materia oscura determinan la morfología final de las galaxias resultantes.

\subsubsection*{Dependencia de parámetros orbitales}

Las colisiones frontales producen las perturbaciones más violentas, mientras que encuentros tangenciales generan interacciones más moderadas.
La orientación relativa de los discos galácticos afecta la formación de barras, ondas espirales y otras estructuras internas.

\subsubsection{Simulaciones emblemáticas en astrofísica}
\subsubsection*{Simulación de las galaxias Antennae}
Este sistema es un ejemplo clásico de colisión avanzada entre dos galaxias espirales. Las simulaciones numéricas han reproducido sus extensas colas de marea y la concentración de gas en sus núcleos, confirmando el impacto de las fuerzas gravitacionales y la interacción del gas.

\subsubsection*{Simulación de la fusión Vía Láctea-Andrómeda}
Simulaciones predictivas muestran que, en aproximadamente 4 mil millones de años, estas dos galaxias espirales colisionarán, formando una galaxia elíptica gigante. Estas simulaciones han proporcionado un marco para comprender cómo las grandes galaxias evolucionan a través de fusiones mayores.

\subsubsection*{Simulaciones cosmológicas a gran escala (Illustris y EAGLE)}
Aunque no se centran exclusivamente en colisiones galácticas individuales, estas simulaciones abordan interacciones en un contexto cosmológico, revelando cómo las fusiones galácticas impulsan la formación de estructuras a gran escala en el universo.

\subsubsection{Limitaciones y desafíos}
Si bien las simulaciones numéricas han sido extremadamente útiles, enfrentan varias limitaciones entre las más destacadas véase.

\begin{enumerate}
    \item Resolución computacional:
Las simulaciones deben equilibrar precisión y tiempo de cálculo. Los modelos más detallados a menudo requieren supercomputadoras y años de procesamiento.
    \item Modelado del gas y formación estelar:
La física del gas y los procesos estelares, como la retroalimentación de supernovas, aún no están completamente comprendidos, lo que introduce incertidumbres en los resultados.
    \item Interacciones de materia oscura:
Aunque se modela como colisión débil, las propiedades exactas de la materia oscura (como autointeracciones potenciales) podrían influir en los resultados y aún están sujetas a investigación.
\end{enumerate}

\section{Capítulo 3: Efectos Extragalácticos}
\subsection{Formación de estructuras a gran escala}

Las colisiones galácticas no solo transforman las galaxias individuales involucradas, sino que también desempeñan un papel crucial en la formación y evolución de estructuras a gran escala en el universo. Desde cúmulos y supercúmulos hasta filamentos cósmicos, las interacciones galácticas contribuyen al entramado jerárquico que define la distribución de materia en el cosmos.

\subsubsection{La conexión entre colisiones galácticas y la estructura cósmica}
El universo visible está estructurado en una red cósmica compuesta por galaxias agrupadas en filamentos y separadas por vastos vacíos. Esta organización es el resultado de la expansión del universo y de los procesos gravitacionales que agrupan la materia a lo largo del tiempo. Las colisiones galácticas son eventos inevitables dentro de esta red y tienen impactos significativos en la evolución de estas estructuras. \\

La estructura jerárquica del universo implica que las galaxias primero se forman en pequeños cúmulos, que luego se agrupan en cúmulos mayores y, finalmente, en supercúmulos. En cada una de estas escalas, las colisiones galácticas son fundamentales para:
\begin{enumerate}
    \item Redistribuir la materia y el momento angular.
    \item Fusionar galaxias más pequeñas en sistemas más grandes.
    \item Contribuir a la evolución de los cúmulos mediante la fusión de sus componentes galácticos.
\end{enumerate}
\subsubsection{Formación de cúmulos galácticos}
\subsubsection*{3.1.1.1 Rol de las colisiones en los cúmulos}
Los cúmulos de galaxias son las mayores estructuras ligadas gravitacionalmente en el universo, con cientos o miles de galaxias interaccionando dentro de un halo compartido de materia oscura. Las colisiones galácticas dentro de estos entornos son frecuentes debido a la densidad extrema de galaxias, lo que resulta en fenómenos como:

\subsubsection*{Acreción de masa}
Las fusiones galácticas incrementan la masa total del cúmulo, que a su vez profundiza el potencial gravitacional del sistema.
\subsubsection*{Interacciones múltiples}
Las galaxias en los cúmulos experimentan interacciones repetidas que pueden despojar a las galaxias de su gas y estrellas exteriores, contribuyendo al medio intracúmulo (ICM). Este gas caliente es detectable en rayos X, y su distribución proporciona información clave sobre las interacciones galácticas.
\subsubsection*{3.1.1.2 Fusión de cúmulos de galaxias}
Cuando dos cúmulos galácticos interactúan, se forman estructuras aún mayores. Ejemplos notables incluyen el "Cúmulo Bala", donde las simulaciones y observaciones sugieren que la materia visible y la materia oscura se comportan de forma distinta durante la colisión. Estos eventos a gran escala son esenciales para comprender la dinámica del universo y para restringir las propiedades de la materia oscura.

\subsubsection{Formación de filamentos cósmicos y vacíos}
\subsubsection*{Colisiones como promotores de los filamentos}
Los filamentos cósmicos son estructuras alargadas que conectan cúmulos de galaxias y forman el armazón de la red cósmica. Las simulaciones cosmológicas muestran que las colisiones y fusiones galácticas dentro de los filamentos contribuyen al flujo de materia hacia los cúmulos, acrecentándolos progresivamente. \\
La redistribución de gas y materia estelar durante las colisiones enriquece los filamentos con elementos químicos y poblaciones estelares difusas.
\subsubsection*{Vacíos cósmicos y la falta de colisiones:}
En contraste, los vacíos cósmicos, que contienen muy pocas galaxias, se caracterizan por una escasez de colisiones galácticas. Esta falta de interacciones refuerza las diferencias evolutivas entre galaxias en los vacíos y aquellas en entornos densos. Las galaxias en los vacíos tienden a ser sistemas menos perturbados, conservando su morfología espiral por más tiempo.

\subsubsection{Efectos en la metalicidad y enriquecimiento químico}
Las colisiones galácticas juegan un papel importante en el enriquecimiento químico del medio intergaláctico y de las estructuras a gran escala.

\subsubsection*{Expulsión de material enriquecido}
Los procesos de marea y los estallidos de formación estelar inducidos por colisiones arrojan gas enriquecido con elementos pesados hacia el medio intergaláctico, contribuyendo a la metalicidad general de los filamentos y cúmulos.
\subsubsection*{Mezcla de poblaciones estelares}
Las fusiones galácticas mezclan estrellas de diferentes generaciones, homogenizando la composición química de las estructuras mayores.
\subsubsection{Implicaciones cosmológicas y materia oscura}
Las colisiones galácticas también son una herramienta clave para estudiar la naturaleza de la materia oscura. Durante las colisiones, se observan fenómenos como.

\subsubsection*{Separación entre materia oscura y materia bariónica}
En sistemas como el "Cúmulo Bala", las observaciones muestran que la materia visible interactúa a través de procesos colisionables (como fricción del gas), mientras que la materia oscura parece ser colisión débil. Esto proporciona límites sobre la naturaleza de la materia oscura, como su sección transversal de interacción.
\subsubsection*{Evolución gravitacional de la red cósmica}
Las simulaciones muestran que los halos de materia oscura de las galaxias dominan la dinámica de los cúmulos y filamentos, influenciando cómo las colisiones redistribuyen la masa en estas estructuras.
\subsubsection{Evidencias observacionales}
Los datos obtenidos mediante telescopios y observatorios espaciales han sido fundamentales para confirmar el impacto de las colisiones galácticas en estructuras a gran escala. Ejemplos destacados incluyen: \\

Levantamientos como el Sloan Digital Sky Survey (SDSS): Este proyecto ha revelado la red cósmica en detalle, confirmando la importancia de las interacciones galácticas en la evolución de cúmulos y filamentos.\\
Observaciones de rayos X y lentes gravitacionales: Estas técnicas han sido clave para estudiar la distribución de gas caliente y materia oscura en cúmulos colisionantes.

\subsection{Influencia en el entorno cósmico (\textit{corrientes de marea, cúmulos de galaxias})}
Las colisiones galácticas generan una variedad de efectos que trascienden a las galaxias individuales involucradas, influyendo de manera significativa en el entorno cósmico. Estos impactos abarcan la formación de corrientes de marea que enriquecen el medio intergaláctico, así como el papel crucial en la evolución de cúmulos galácticos y su dinámica interna.

\subsubsection{Corrientes de marea: puentes y colas estelares}
\subsubsection*{Origen de las corrientes de marea}
Las fuerzas de marea durante una colisión galáctica son el resultado de las diferencias gravitacionales entre las distintas regiones de una galaxia cuando interactúa con otra. Estas fuerzas pueden arrancar estrellas y gas de las galaxias involucradas, creando estructuras alargadas conocidas como puentes y colas de marea.

\subsubsection*{Características principales}

Los puentes de marea se forman entre dos galaxias durante la colisión, conectándolas con un flujo de material. Un ejemplo icónico es el puente entre las galaxias Antennae, donde el material eyectado es visible en longitudes de onda ópticas e infrarrojas. Las colas de marea Son estructuras alargadas que se extienden hacia el medio intergaláctico, compuestas principalmente de estrellas y gas eyectado. La cola de la galaxia M82, por ejemplo, muestra evidencia de enriquecimiento químico debido a estallidos de formación estelar.

\subsubsection*{Importancia astrofísica}

Las corrientes de marea redistribuyen masa y momento angular, afectando la evolución de las galaxias involucradas. Enriquecen el medio intergaláctico con gas y estrellas eyectadas, alterando su composición química y densidad. Pueden dar lugar a la formación de nuevas galaxias enanas en los nodos de mayor densidad del material eyectado, un fenómeno conocido como galaxias de marea.
\subsubsection{Evolución de los cúmulos galácticos}
\subsubsection*{Interacciones dinámicas en los cúmulos}
En los cúmulos de galaxias, donde la densidad de galaxias es alta, las colisiones galácticas son comunes y contribuyen significativamente a la evolución dinámica del cúmulo. \\

Las interacciones gravitacionales entre galaxias producen intercambios de energía, lo que puede alterar sus trayectorias orbitales dentro del cúmulo.
Las colisiones repetidas desgastan las estructuras galácticas, despojándolas de su gas externo mediante un proceso conocido como stripping de marea.
\subsubsection*{Efectos en el medio intracúmulo (ICM)}
El medio intracúmulo es el gas caliente que permea los cúmulos de galaxias, detectable en rayos X. Las colisiones galácticas contribuyen al ICM de varias maneras. \\

\noindent \textit{Expulsión de gas galáctico:}
Durante las colisiones, el gas interestelar es eyectado y se mezcla con el ICM, enriqueciendo este medio con elementos pesados. \\
\textit{Ondas de choque:}
Las colisiones generan ondas de choque en el gas intracúmulo, calentándolo aún más y redistribuyendo la energía térmica en el cúmulo.
\subsubsection*{Formación de cúmulos relajados y no relajados}

Los cúmulos relajados, como Virgo, muestran un núcleo denso y estable donde las colisiones han permitido que las galaxias pierdan energía y se acumulen en el centro.
Los cúmulos no relajados, como el "Cúmulo Bala", exhiben distribuciones irregulares de materia y gas debido a colisiones recientes entre cúmulos menores o grandes galaxias.
\subsubsection{Redistribución de materia y formación de estructuras}
\subsubsection*{Papel en la red cósmica}

Las colisiones galácticas redistribuyen masa hacia los filamentos cósmicos y los cúmulos, contribuyendo al crecimiento jerárquico de la estructura cósmica.
Los remanentes de fusiones galácticas a menudo se convierten en galaxias centrales masivas (cD), que dominan los cúmulos galácticos en términos de masa y luminosidad.
\subsubsection*{Influencia gravitacional a gran escala}
Las galaxias involucradas en colisiones pueden actuar como lentes gravitacionales, distorsionando la luz de objetos de fondo. Este efecto permite a los astrónomos mapear la distribución de masa, incluida la materia oscura, en los cúmulos.

\subsubsection{Evidencias observacionales}
\subsubsection*{Observaciones de rayos X}

Los telescopios de rayos X, como Chandra y XMM-Newton, han proporcionado imágenes detalladas del ICM, mostrando ondas de choque y estructuras asociadas con colisiones galácticas y cúmulos en interacción.
\subsubsection*{Ejemplos destacados}

\textit{Cúmulo Bala:} \\
Observado como un ejemplo paradigmático de colisión entre dos cúmulos, donde la separación entre la materia visible y la materia oscura proporciona evidencia de las propiedades de la materia oscura. \\
\textit{Sistema de galaxias Antennae:} \\
Este par en colisión presenta espectaculares corrientes de marea y regiones de formación estelar intensificada.
\subsubsection{Implicaciones cosmológicas}
\subsubsection*{Formación de galaxias y cúmulos masivos}
Las colisiones galácticas dentro de los cúmulos contribuyen al crecimiento de las galaxias centrales y al enriquecimiento de su medio circundante.

\subsubsection*{Relación con la materia oscura}
El comportamiento de las colisiones, especialmente en cúmulos como el "Cúmulo Bala", brinda restricciones observacionales sobre la naturaleza y propiedades de la materia oscura, como su colisión débil y distribución gravitacional.

\subsection{Relación con la evolución de las galaxias anfitrionas}
Las colisiones galácticas desempeñan un papel transformador en la evolución de las galaxias anfitrionas. Estos eventos generan cambios significativos en la morfología, la dinámica interna, el contenido de gas y estrellas, y la actividad de formación estelar y nuclear. A continuación, se describen los principales impactos que las colisiones tienen en las galaxias involucradas.

\subsubsection{Transformaciones morfológicas}
\subsubsection*{Cambios en la estructura galáctica}

Las interacciones gravitacionales durante las colisiones deforman las galaxias anfitrionas, generando estructuras como colas de marea, anillos estelares y puentes de materia. Las colisiones menores (o fusiones menores) pueden alterar sutilmente las morfologías de galaxias espirales, mientras que las colisiones mayores (o fusiones mayores) a menudo resultan en la transformación de galaxias espirales en elípticas.
\subsubsection*{Formación de galaxias elípticas}

Las fusiones mayores tienden a destruir los discos galácticos y redistribuir las estrellas en configuraciones esferoidales, formando galaxias elípticas.
Este proceso, observado en simulaciones numéricas y en sistemas reales como las galaxias Antennae, es clave para entender la evolución de las galaxias masivas en el universo.
\subsubsection{Actividad de formación estelar}
\subsubsection*{Estallidos de formación estelar inducidos}
Las colisiones comprimen el gas interestelar en las galaxias anfitrionas, disparando episodios intensos de formación estelar conocidos como starbursts. Durante estas fases, las tasas de formación estelar pueden aumentar hasta 100 veces respecto a niveles normales.
Ejemplo: Las galaxias en interacción Arp 220 muestran tasas de formación estelar extremas, alimentadas por gas comprimido durante su colisión.
\subsubsection*{Distribución de regiones de formación estelar}

La formación estelar puede concentrarse en los núcleos galácticos, produciendo núcleos luminosos infrarrojos.
También puede extenderse a lo largo de las colas de marea y los puentes, generando poblaciones estelares en estas regiones.
\subsubsection*{Consecuencias a largo plazo}

Los estallidos de formación estelar agotan rápidamente el gas disponible, dejando galaxias con bajas reservas de combustible para futuras generaciones estelares.
Esto contribuye a la transición de galaxias espirales a sistemas más pasivos y dominados por estrellas antiguas.
\subsubsection{Actividad nuclear: galaxias activas y cuásares}
\subsubsection*{Alimentación de agujeros negros supermasivos}

Durante las colisiones, las fuerzas gravitacionales canalizan gas hacia los núcleos galácticos, alimentando los agujeros negros supermasivos (AGN, por sus siglas en inglés). Este gas puede activar núcleos galácticos activos, produciendo emisión en múltiples longitudes de onda, desde rayos X hasta ondas de radio.
\subsubsection*{Transición a cuásares}

Las fusiones mayores pueden generar cuásares, las fases más luminosas de los AGN, donde el agujero negro central libera grandes cantidades de energía debido al acrecimiento de material. Ejemplo: Los cuásares observados en sistemas de fusión como NGC 6240 evidencian la relación entre las colisiones y la actividad nuclear extrema.
\subsubsection*{Relación con la evolución de las galaxias anfitrionas}

La retroalimentación del AGN, en forma de vientos galácticos y jets relativistas, puede expulsar gas del núcleo y suprimir la formación estelar, contribuyendo al agotamiento del gas frío en las galaxias anfitrionas.
\subsubsection{Redistribución de masa y momento angular}
\subsubsection*{Pérdida de momento angular}

Durante las interacciones, el gas y las estrellas pierden momento angular, redistribuyéndose hacia el núcleo galáctico o hacia el entorno intergaláctico.
Esta pérdida de momento angular es clave para alimentar tanto los AGN como la formación estelar centralizada.
\subsubsection*{Formación de halos estelares y gas difuso}

Las colisiones expulsan material hacia los alrededores de las galaxias anfitrionas, formando halos de estrellas y gas difuso. Este material contribuye a la metalicidad del entorno intergaláctico y puede ser reincorporado en etapas posteriores de evolución galáctica.
\subsubsection{Evidencias observacionales y simulaciones numéricas}
\subsubsection*{Sistemas en colisión}

Las galaxias Antennae son un ejemplo paradigmático de las transformaciones inducidas por colisiones, con tasas de formación estelar intensas y estructuras de marea prominentes.
El sistema Arp 220 destaca por su actividad infrarroja extrema y su núcleo doble, evidencia de una fusión mayor en curso.
\subsubsection*{Apoyo de simulaciones}

Las simulaciones hidrodinámicas y de N-cuerpos reproducen con precisión los efectos observados, mostrando cómo las colisiones alteran las morfologías, activan los núcleos galácticos y generan corrientes de marea. Estas simulaciones son fundamentales para estudiar colisiones en diferentes etapas y comprender los procesos subyacentes.
\subsubsection{Implicaciones para la evolución cósmica}
\subsubsection*{Rol en la historia del universo}

Las colisiones galácticas fueron más comunes en el universo temprano, cuando las galaxias estaban más próximas debido a la menor expansión cósmica.
Estos eventos contribuyeron al ensamblaje jerárquico de galaxias masivas y cúmulos galácticos, definiendo la estructura actual del universo.
\subsubsection{Conexión con la materia oscura}

Las propiedades dinámicas de las galaxias durante las colisiones dependen de la distribución de materia oscura en sus halos.
Esto implica que las colisiones no solo afectan a las galaxias anfitrionas, sino que también ofrecen información valiosa sobre la naturaleza de la materia oscura.

\section{Capítulo 4: Rol de la Materia Oscura}
\subsection{Dinámica gravitacional en sistemas dominados por materia oscura}
La materia oscura desempeña un papel crucial en la dinámica gravitacional de los sistemas galácticos y extragalácticos. Aunque no interactúa directamente con la radiación electromagnética, su influencia gravitacional domina el comportamiento de las galaxias, los cúmulos y las estructuras a gran escala. Este capítulo explora cómo la presencia de materia oscura afecta la dinámica de los sistemas galácticos involucrados en colisiones, así como su rol en la estabilidad y evolución de estos sistemas.

\subsubsection{Distribución de materia oscura en halos galácticos}
\subsubsection*{Composición y estructura}

La mayoría de las galaxias están rodeadas por un halo de materia oscura, una estructura esférica que contiene la mayor parte de la masa galáctica. Estos halos no son visibles, pero su presencia se infiere a partir de efectos gravitacionales, como las curvas de rotación galáctica planas observadas en muchas espirales.
\subsubsection*{Importancia dinámica}

Los halos de materia oscura estabilizan las galaxias, actuando como un "pegamento gravitacional" que mantiene unidas a las estrellas y al gas, incluso durante colisiones violentas. En los sistemas de colisión, los halos determinan las trayectorias orbitales y la eficiencia con la que se transfiere energía cinética entre las galaxias.
\subsubsection*{Relación con la densidad de la materia}

Las simulaciones de formación galáctica sugieren que los halos de materia oscura siguen una distribución conocida como perfil de Navarro-Frenk-White (NFW), con densidad creciente hacia el centro y decreciente en las regiones externas.
\subsubsection{Efectos de la materia oscura en las colisiones galácticas}
\subsubsection*{Modulación de la dinámica de fusión}

La presencia de materia oscura altera el comportamiento de las colisiones al incrementar la masa gravitacional efectiva de las galaxias.
Los halos permiten que las galaxias interactúen gravitacionalmente incluso a grandes distancias, facilitando la captura mutua y la eventual fusión.
\subsubsection*{Transferencia de momento angular y energía}

Durante una colisión, el halo de materia oscura absorbe y redistribuye el momento angular y la energía cinética, amortiguando los efectos sobre las componentes visibles (estrellas y gas).
Este efecto permite que las galaxias sobrevivan a múltiples interacciones sin desintegrarse completamente.
\subsubsection*{Formación de estructuras transitorias}

Las fusiones entre galaxias dominadas por materia oscura a menudo producen estructuras temporales, como puentes de marea y anillos, antes de alcanzar un estado de equilibrio. La materia oscura estabiliza estas estructuras al proporcionar un campo gravitacional subyacente robusto.
\subsubsection{Influencia en el enriquecimiento del entorno galáctico}
\subsubsection*{Expulsión de gas galáctico}

La materia oscura influye en la eficiencia con la que el gas es expulsado durante las colisiones.
Halos más masivos pueden retener material eyectado, mientras que halos más pequeños permiten que el gas escape hacia el medio intergaláctico, contribuyendo al enriquecimiento químico.
\subsubsection*{Generación de ondas de choque}

En sistemas dominados por materia oscura, las colisiones galácticas generan ondas de choque en el gas intragaláctico. Estas ondas redistribuyen energía térmica, enriquecen el entorno y afectan la formación estelar en las regiones circundantes.
\subsubsection{Evidencia observacional: Materia oscura y colisiones}
\subsubsection*{Separación de materia visible y oscura}

Un ejemplo clave es el Cúmulo Bala, donde la materia oscura y el gas caliente muestran una clara separación debido a sus diferentes formas de interacción.
Observaciones de lentes gravitacionales en este cúmulo revelan que la mayor parte de la masa se encuentra en la materia oscura, mientras que el gas visible experimenta interacciones viscosas.
\subsubsection*{Lentes gravitacionales en galaxias en interacción}

Los sistemas de colisión actúan como lentes gravitacionales fuertes o débiles, deformando la luz de objetos de fondo. Esto permite mapear la distribución de materia oscura en los sistemas anfitriones.
\subsubsection{Importancia en simulaciones numéricas}
\subsubsection*{Inclusión de materia oscura en modelos}

Las simulaciones modernas de colisiones galácticas incorporan materia oscura como componente fundamental, modelada mediante partículas de N-cuerpos con interacciones gravitacionales. Estas simulaciones han reproducido estructuras observadas, como las colas de marea y las galaxias elípticas formadas tras fusiones mayores.
\subsubsection*{Impacto en la formación de galaxias}

Los modelos sugieren que la materia oscura regula la formación estelar durante las colisiones al controlar la cantidad de gas disponible y su compresión. La interacción entre el gas visible y la materia oscura también explica la aparición de galaxias enanas satélites en las proximidades de sistemas de colisión.
\subsubsection{Rol cosmológico: escalas mayores}
\subsubsection*{Ensamblaje de estructuras a gran escala}

La materia oscura guía la formación jerárquica de estructuras, desde galaxias individuales hasta cúmulos galácticos y filamentos cósmicos. Durante colisiones, las galaxias afectadas tienden a seguir la red gravitacional dictada por los halos de materia oscura, contribuyendo al crecimiento de cúmulos masivos.
\subsubsection*{Implicaciones para la teoría de la materia oscura}

Las observaciones de colisiones proporcionan restricciones clave sobre la naturaleza de la materia oscura, como su falta de auto-interacción significativa. Experimentos como los realizados en el Cúmulo Bala han descartado muchas teorías alternativas, favoreciendo modelos de materia oscura fría (CDM).

\subsection{Impacto en las trayectorias y resultados de las colisiones}
La materia oscura juega un papel crucial en la configuración de las trayectorias de las galaxias en interacción y en los resultados de sus colisiones. Al ser la componente dominante en términos de masa gravitacional, los halos de materia oscura influyen directamente en la forma en que las galaxias interactúan, las estructuras que generan, y su evolución tras la colisión. Este apartado aborda cómo la materia oscura afecta estos aspectos clave.

\subsubsection{Modulación de las trayectorias orbitales}
\subsubsection*{Aumento de la atracción gravitacional}

Los halos de materia oscura amplifican la fuerza gravitacional entre las galaxias en interacción, incrementando las probabilidades de capturas gravitacionales y encuentros cercanos. Esto resulta en trayectorias más cerradas y recurrentes, facilitando la eventual fusión de las galaxias.
\subsubsection*{Desviaciones respecto a trayectorias keplerianas}

Las trayectorias de las galaxias no siguen órbitas puramente keplerianas debido a la distribución extendida de masa en los halos. Las interacciones gravitacionales entre halos generan fenómenos como órbitas elongadas y cambios de dirección abruptos durante acercamientos cercanos.
\subsubsection*{Decaimiento orbital}

La fricción dinámica, causada por la interacción entre los halos de materia oscura y las partículas estelares, provoca un decaimiento gradual de las órbitas de las galaxias en interacción, acelerando su coalescencia.
\subsubsection{Redistribución de masa y momento angular}
\subsubsection*{Transferencia de momento angular}

Durante las colisiones, los halos de materia oscura absorben y redistribuyen el momento angular. Este proceso es clave para transformar órbitas iniciales de alta excentricidad en configuraciones estables post-colisión. La pérdida de momento angular permite que el gas y las estrellas se acumulen en los núcleos galácticos, alimentando procesos como formación estelar intensa o actividad de núcleos galácticos activos (AGN).
\subsubsection*{Expulsión de materia visible}

La interacción entre las componentes visibles (gas y estrellas) y la materia oscura puede causar la eyección de material a grandes distancias, formando estructuras como colas de marea y corrientes estelares. Estas estructuras son más pronunciadas en sistemas donde los halos son asimétricos o tienen densidades centrales bajas.
\subsubsection{Influencia en los resultados de las colisiones}
\subsubsection*{Formación de nuevas estructuras galácticas}

La presencia de halos masivos de materia oscura favorece la fusión completa de las galaxias, generando sistemas resultantes como galaxias elípticas masivas.
En colisiones menores, los halos permiten la supervivencia de las galaxias satélites, que pueden transformarse en sistemas enanos dominados por materia oscura.
\subsubsection*{Generación de subestructuras}

Las simulaciones han demostrado que las interacciones gravitacionales entre los halos producen corrientes estelares y cúmulos satelitales. Estas estructuras son útiles para rastrear la distribución de materia oscura en los sistemas anfitriones.
\subsubsection*{Influencia en la actividad nuclear}

El papel estabilizador de la materia oscura en las colisiones permite la acumulación de gas en los núcleos galácticos, lo que puede desencadenar actividad AGN. Este proceso tiene importantes implicaciones para la evolución de las galaxias involucradas.
\subsubsection{Evidencias observacionales y simulaciones numéricas}
\subsubsection*{Sistemas como el Cúmulo Bala}

En colisiones a gran escala, la materia oscura afecta la dinámica de los cúmulos galácticos, como se evidencia en el Cúmulo Bala, donde los núcleos de materia oscura permanecen separados de las componentes visibles tras la colisión.
\subsubsection*{Simulaciones avanzadas}

Simulaciones hidrodinámicas con halos de materia oscura han demostrado cómo la distribución de masa afecta las trayectorias y los resultados de las colisiones. Por ejemplo, los modelos de colisiones espirales-elípticas reproducen las transformaciones morfológicas observadas y las tasas de formación estelar inducida.
\subsubsection{Dependencia de las propiedades del halo de materia oscura}
\subsubsection*{Tamaño y masa del halo}

Halos más masivos generan interacciones más fuertes, aumentando las probabilidades de colisión y la severidad de los eventos.
Halos menos masivos, comunes en galaxias enanas, permiten interacciones menos disruptivas, preservando la morfología de las galaxias involucradas.
\subsubsection*{Concentración del halo}

Halos con mayor concentración central inducen fusión más rápida y eficiente, mientras que halos extendidos y difusos permiten interacciones prolongadas y multietapas.
\subsubsection{Conexión con la cosmología jerárquica}
\subsubsection*{Formación de galaxias y cúmulos}

En el marco de la cosmología de materia oscura fría (CDM), las colisiones galácticas son etapas clave en el ensamblaje jerárquico de estructuras.
La presencia de halos permite que las galaxias en interacción se integren en cúmulos galácticos sin dispersarse en el medio intergaláctico.
\subsubsection*{Implicaciones para el modelo cosmológico}

Las observaciones de trayectorias y resultados de colisiones restringen parámetros clave del modelo CDM, como la auto-interacción y la distribución inicial de los halos.

\section{*Capítulo Extra: Ejemplos Actuales y $\mathfrak{Milkomeda}$}
Las colisiones galácticas no son solo fenómenos teóricos o históricos; en el cosmos actual y futuro próximo, se observan ejemplos emblemáticos de galaxias en interacción que revelan detalles cruciales sobre estos procesos. Este capítulo examina dos casos específicos: las galaxias Antena, un ejemplo icónico de colisión en curso, y la futura colisión entre la Vía Láctea y Andrómeda, que dará lugar a la galaxia conocida como "Milkomeda". \\

\subsection{Las Galaxias Antena (NGC 4038 y NGC 4039): Colisión en curso}
\subsubsection*{Descripción general}

Las galaxias Antena, situadas a aproximadamente 45 millones de años luz de la Tierra, en la constelación de Corvus, representan un ejemplo arquetípico de fusión galáctica avanzada.
Este sistema debe su nombre a las estructuras de "antena" que se extienden hacia afuera, formadas por corrientes de marea inducidas por la interacción gravitacional entre ambas galaxias.
\subsubsection*{Estructuras clave observadas}

\textit{Núcleos galácticos en fusión:} Ambas galaxias han comenzado a fusionarse, aunque sus núcleos aún son distinguibles. En el futuro, se espera que formen un único núcleo galáctico masivo. \\
\textit{Colas de marea:} Estas estructuras extendidas están formadas por gas y estrellas arrancadas de las galaxias debido a fuerzas gravitacionales intensas. \\
\textit{Regiones de formación estelar:} El gas comprimido por la interacción gravitacional ha dado lugar a regiones de intensa formación estelar, evidentes en longitudes de onda óptica e infrarroja.
\subsubsection*{Relevancia científica}

Las galaxias Antena son un laboratorio natural para estudiar los efectos de las colisiones galácticas, como la redistribución de masa, el aumento en la formación estelar y la evolución morfológica. Observaciones en múltiples longitudes de onda, desde rayos X hasta radio, han proporcionado información sobre la distribución de gas, estrellas y materia oscura en este sistema.
\subsection{Colisión futura entre la Vía Láctea y Andrómeda (Milkomeda)}
\subsubsection*{Contexto y cronología del evento}

La galaxia Andrómeda (M31), situada a unos 2.5 millones de años luz de distancia, se está acercando a la Vía Láctea a una velocidad de aproximadamente 110 km/s debido a su atracción gravitacional mutua y la influencia de la materia oscura en sus halos. \\
En unos 4 a 5 mil millones de años, ambas galaxias comenzarán a interactuar, entrando en un proceso que culminará en su fusión para formar una galaxia elíptica gigante conocida como Milkomeda.
\subsubsection*{Fases de la colisión}

\textit{Primera aproximación:} Las galaxias se acercarán lo suficiente como para que sus halos interactúen, generando corrientes de marea y un aumento en la formación estelar. \\
\textit{Fase de interacción repetida:} A lo largo de miles de millones de años, las galaxias pasarán cerca repetidamente, perdiendo momento angular y redistribuyendo masa en cada encuentro. \\
\textit{Fusión final:} Eventualmente, los núcleos de ambas galaxias se fusionarán, dando lugar a una estructura elíptica masiva.
\subsubsection*{Efectos en la Vía Láctea}

\textit{Redistribución de las estrellas y gas:} Muchas estrellas, incluido el Sol, serán desplazadas de sus órbitas actuales y podrán ser expulsadas al halo galáctico o eyectadas al espacio intergaláctico. \\
\textit{Disminución de la formación estelar:} La fusión agotará la mayor parte del gas disponible, lo que llevará a una disminución gradual en la formación de nuevas estrellas. \\
\textit{Actividad nuclear transitoria:} El gas acumulado en el núcleo resultante podría desencadenar un núcleo galáctico activo (AGN) temporal.
\subsubsection*{Importancia cosmológica}

La futura colisión entre la Vía Láctea y Andrómeda es un ejemplo de la evolución jerárquica del universo, donde las galaxias pequeñas y medianas se combinan para formar sistemas mayores. Este evento proporcionará información clave sobre el papel de la materia oscura en la dinámica a gran escala y la evolución de las galaxias.
\subsection*{Similitudes y diferencias entre Antena y Milkomeda}

Ambos eventos destacan el rol central de la materia oscura en guiar las trayectorias de las galaxias en colisión. En ambos casos, las interacciones provocan procesos de marea, redistribución de masa y formación de nuevas estructuras; las galaxias Antena están en una fase avanzada de fusión, mientras que la colisión Vía Láctea-Andrómeda es un evento futuro en sus etapas iniciales. Las Antena son más pequeñas y menos masivas en comparación con Milkomeda, lo que implica escalas temporales y dinámicas diferentes.

\subsection*{Conclusiones}

A lo largo de este trabajo, se han explorado los principales aspectos relacionados con la fenomenología de las colisiones galácticas, abarcando sus fundamentos teóricos, mecanismos, efectos extragalácticos, y el papel crucial de la materia oscura, además de ilustrar el fenómeno mediante ejemplos observacionales y eventos futuros. Cada capítulo ha contribuido a construir una perspectiva integral sobre la relevancia de estas interacciones en el contexto de la evolución cósmica y su impacto en el modelo actual de formación y estructura del universo. \\

En el Capítulo 1, se estableció el marco conceptual que sustenta el estudio de las colisiones galácticas. Desde la dinámica galáctica básica hasta las diferentes formas de interacción —como fusiones y encuentros cercanos—, se destacó cómo estos procesos constituyen eventos transformadores en la evolución de las galaxias. Asimismo, se introdujo el concepto de materia oscura, destacando su presencia dominante en las galaxias y su influencia en sus trayectorias y estructuras, un tema que se desarrolló más adelante en el trabajo. \\

El Capítulo 2 profundizó en los mecanismos físicos que impulsan las colisiones galácticas. Desde las interacciones gravitacionales que generan fuerzas de marea hasta la redistribución de masa estelar, se identificaron los procesos fundamentales que alteran tanto la estructura interna como la dinámica global de las galaxias involucradas. Además, el papel de las simulaciones numéricas como herramienta esencial para comprender estas interacciones quedó evidenciado, proporcionando una visión detallada de los posibles resultados de estos eventos y de las implicaciones a largo plazo para las estructuras resultantes. \\

El Capítulo 3 abordó los efectos extragalácticos derivados de las colisiones, subrayando su papel en la formación de estructuras a gran escala, como cúmulos de galaxias y filamentos cósmicos. También se discutió la influencia de las interacciones en el entorno cósmico, como la creación de corrientes de marea y subestructuras satelitales, y su relación con la evolución de las galaxias anfitrionas, revelando cómo estos eventos están interconectados con la jerarquía cosmológica del universo. \\

El Capítulo 4 centró la atención en el rol fundamental de la materia oscura, cuya presencia no solo determina las trayectorias y resultados de las colisiones, sino que también estabiliza y guía la formación de estructuras a partir de los eventos de interacción. Las evidencias observacionales y las simulaciones numéricas demuestran que los halos de materia oscura son responsables de procesos críticos como el decaimiento orbital, la acumulación de gas en los núcleos galácticos y la formación de galaxias masivas post-colisión, consolidando su papel como pieza clave en la dinámica universal.  \\

Finalmente, en el Capítulo 5, se revisaron ejemplos concretos y proyecciones futuras de colisiones galácticas. Las galaxias Antena, como caso emblemático de interacción en curso, han permitido estudiar los efectos de las fuerzas gravitacionales y de marea en un sistema observable, mientras que la colisión futura entre la Vía Láctea y Andrómeda (Milkomeda) ejemplifica la naturaleza inevitable y cíclica de estos eventos en el marco de la evolución jerárquica del universo. Ambos casos refuerzan la idea de que las colisiones galácticas no son eventos aislados, sino procesos continuos que moldean la estructura del cosmos.  \\

Las colisiones galácticas representan un fenómeno multifacético cuya comprensión requiere la integración de teorías físicas, observaciones detalladas y simulaciones avanzadas. Estos eventos son motores fundamentales de la evolución galáctica, desempeñando un papel central en la redistribución de masa, el desencadenamiento de procesos estelares y la formación de estructuras a gran escala. La materia oscura, como protagonista invisible, influye de manera decisiva en cada etapa de las colisiones, desde las trayectorias iniciales hasta las configuraciones finales. Este trabajo no solo resalta la importancia de estas interacciones en la dinámica universal, sino que también contribuye a consolidar nuestra comprensión del papel de la materia oscura en el cosmos, un área de investigación que sigue siendo esencial para desentrañar los misterios del universo.  \\

\section*{Agradecimientos}

A la profesora la doctora Margarita Eugenia del Socorro Rosado Solís por habernos introducido a la astronomía y tópicos de astrofísica a lo largo del semestre 2025-1, siendo este nuestro comienzo de nuestra trayectoria como futuros científicos apasionados por los cuerpos celestes y desvelar los misterios que este nos ofrece. Así como al maestro en astronomía Isidro Ramírez Ballinas por contribuir a darnos una pequeña segunda perspectiva del estudio de esta ciencia.

\section*{Bibliografía}

\begin{itemize}
\item Binney, J., & Tremaine, S. (2008). *Galactic dynamics: Second edition*. Princeton University Press.
\item Mo, H. J., van den Bosch, F. C., & White, S. D. M. (2010). *Galaxy formation and evolution*. Cambridge University Press.
\item Peebles, P. J. E. (1993). *Principles of physical cosmology*. Princeton University Press.
\item Sparke, L. S., & Gallagher, J. S. (2007). *Galaxies in the universe: An introduction* (2nd ed.). Cambridge University Press.
\item Longair, M. S. (2008). *Galaxy formation*. Springer.
\item Fall, S. M., & Efstathiou, G. (1980). *The formation and evolution of galaxies*. Cambridge University Press.
\item Frenk, C. S., & White, S. D. M. (2012). *The formation of cosmic structures*. Oxford University Press.
\item Blanton, M. R., & Moustakas, J. (2009). *Cosmology with large-scale surveys*. Springer.
\item Shapiro, P. R., & Moore, B. (2004). *The physics of galaxies and cosmology*. Springer.
\item Liddle, A. R., & Lyth, D. H. (2000). *Cosmological inflation and large-scale structure*. Cambridge University Press.
\item Bertone, G., Hooper, D., & Silk, J. (2005). *Particle dark matter: Observations, models and searches*. Cambridge University Press.
\item Kolb, E. W., & Turner, M. S. (1990). *The early universe*. Addison-Wesley.
\item Ryden, B. (2003). *Introduction to cosmology*. Pearson.
\item Schneider, P. (2006). *Extragalactic astronomy and cosmology: An introduction*. Springer.
\item Turner, M. S. (1997). *The dark side of the universe: The study of dark matter and dark energy*. Princeton University Press.
\item Carroll, S. M. (2004). *Spacetime and geometry: An introduction to general relativity*. Addison-Wesley.
\item Dodelson, S. (2003). *Modern cosmology*. Academic Press.
\item Viana, P. T. P., & Liddle, A. R. (1996). *The formation of the universe*. Cambridge University Press.
\item Bergström, L., & Goobar, A. (2006). *Cosmology and particle astrophysics*. Springer.
\item Zwicky, F. (1937). *On the mass of nebulae and clusters of nebulae*. Proceedings of the National Academy of Sciences.
\end{itemize}

\end{document}